\input{../tex-inputs/latex-leseansicht-vorspann}

\section[Arthur Schnitzler an Richard Beer-Hofmann, 14. 12. 1896]{L00628 Arthur Schnitzler an Richard Beer-Hofmann, 14. 12. 1896}
\nopagebreak\mylabel{L00628v}
\rehead{ }\normalsize\beginnumbering\briefempfaengerindex{Beer-Hofmann, Richard@\textsc{Beer-Hofmann, Richard}!zzzSchnitzler, Arthur@\emph{von Arthur Schnitzler}!1896-12-141@{14. 12. 1896}|(be}
\toendnotes[C]{\smallbreak\pagebreak[2]}
\correspDesc{Versand  durch Arthur Schnitzler am 14. 12. 1896 in Frankgasse 1
\newline{}Übermittlung  am 14. 12. 1896 in Wien 1/1
\newline{}Erhalt  durch Richard Beer-Hofmann im Zeitraum [14. 12. 1896 – 18. 12. 1896?] in Wollzeile 15 (»Berthahof«)}\toendnotes[C]{\smallbreak}
\Standort{YCGL, MSS 31.}
\physDesc{Postkarte, 112 Zeichen
\newline{}Handschrift: Bleistift, deutsche Kurrent
\newline{}Versand: Stempel: »\nobreak{}\oindex{I., Innere Stadt@\textbf{I., Innere Stadt}|pwk}Wien 1/1, 14. 12. 96, 5–6 V\nobreak{}«.  }\toendnotes[C]{\smallbreak}\pstart{}{\pb}Herrn \textsc{Dr. Rich.
                     Beer-Hofmann}\pend{}\pstart{}Wien\oindex{Wien@\textbf{Wien}|pw}.\pend{}\pstart{}\textsc{I. Wollzeile 15\oindex{Wien@\textbf{Wien}!I., Innere Stadt@\textbf{I., Innere Stadt}!Wollzeile 15 (»Berthahof«)@\textbf{Wollzeile 15 (»Berthahof«)}|pw}}\pend{}{\bigskip}\vspace{1em}
\pstart
           \noindent{}{\pb}Lieber Richard, ich hoffe Sie ko{\geminationm}en
                \label{K_L00628-1v}\edtext{morgen\eventindex{Frankgasse 1@\textbf{Frankgasse 1}!Private Lesung des Prologs zu Mimi, 15.12.1896@Private Lesung des Prologs zu Mimi, 15.12.1896|pwv}}{\lemma{\textnormal{\emph{morgen}}}\Cendnote{\textnormal{Vgl. Private Lesung des Prologs zu Mimi, 15.12.1896. In: A. S.: \emph{Kulturveranstaltungen}, \url{https://schnitzler-kultur.acdh.oeaw.ac.at/pmb335141.html}.}}}\label{K_L00628-1} beſti{\geminationm}t\pend
           \pstart Herzlich Ihr \spacefill\mbox{A. S}\pend{}\selectlanguage{ngerman}\endnumbering\briefempfaengerindex{Beer-Hofmann, Richard@\textsc{Beer-Hofmann, Richard}!zzzSchnitzler, Arthur@\emph{von Arthur Schnitzler}!1896-12-141@{14. 12. 1896}|)be}\mylabel{L00628h}  \newcommand{\dateiname}{L00628}\newcommand{\titel}{Arthur Schnitzler an Richard Beer-Hofmann, 14. 12. 1896}\newcommand{\editorInnen}{Martin Anton Müller und Gerd-Hermann Susen}\input{../tex-inputs/latex-leseansicht-abspann}
